\section{Inclusion Criteria}\label{sec:inclusion-criteria}

For a patient to be included in formal analyses, they needed to have enough \gls{eeg} data and seizures, so the models had enough data to learn patterns for seizure forecasting and for statistical validity of the analysis.
All of the following criteria had to be met:
\begin{itemize}
    \item \textbf{Complete ULTRA2 data availability:} for ULTRA2, only patients who had completed the study (and therefore had their full outpatient recording available) were considered.
    \item \textbf{Sufficient seizure count:} at least $10$ \emph{valid} seizures.
    \item \textbf{Sufficient recording coverage:} recording coverage $\geq\SI{60}{\percent}$.
\end{itemize}

A seizure was considered \emph{valid} if no other seizure occurred within its preictal interval (\cref{subsubsec:interval-types}), i.e., it occurred at least $\sim$\SI{65}{\min} after the previous seizure.

Recording coverage quantifies how much of the total monitoring timespan (from the start of the first recording to the end of the last) contains recorded data:
\[
\text{recording coverage} = \frac{\text{total recorded duration}}{\text{total monitoring timespan}}.
\]

\Cref{tab:patient-inclusion} summarizes which patients met the criteria.

\begin{table}[hptb]
    \centering
    \caption[Patient Inclusion]{Patient Inclusion Criteria.\\
    Next to the valid seizures and recording coverage, it is denoted whether the corresponding criteria were met ($\checkmark$) or not ($\times$). ULTRA2 patients who had not completed the study are omitted.}
    \label{tab:patient-inclusion}
    \setlength{\tabcolsep}{9pt}
    \renewcommand{\arraystretch}{1.15}
    \rowcolors{2}{gray!10}{white} % Alternating row colors
    \begin{tabular}{|l r|r r r|r r r r|}
\hline
\rowcolor{gray!25} \multicolumn{2}{|c|}{Patient} & \multicolumn{3}{c|}{Seizures} & \multicolumn{4}{c|}{Timespan} \\
\rowcolor{gray!25} ID & Incl. & Total & Valid & Met & Monitoring & Recorded & Coverage & Met \\
\hline
C01 & $\checkmark$ & 92 & 58 & $\checkmark$ & 494 & 388 & \SI{78}{\percent} & $\checkmark$ \\
C02 & $\checkmark$ & 54 & 53 & $\checkmark$ & 537 & 487 & \SI{91}{\percent} & $\checkmark$ \\
C03 & $\checkmark$ & 63 & 63 & $\checkmark$ & 415 & 394 & \SI{95}{\percent} & $\checkmark$ \\
U01 & $\checkmark$ & 259 & 195 & $\checkmark$ & 401 & 381 & \SI{95}{\percent} & $\checkmark$ \\
U02 & $\times$ & 2 & 2 & $\times$ & 293 & 102 & \SI{35}{\percent} & $\times$ \\
U03 & $\checkmark$ & 59 & 52 & $\checkmark$ & 394 & 295 & \SI{75}{\percent} & $\checkmark$ \\
U04 & $\checkmark$ & 64 & 62 & $\checkmark$ & 197 & 180 & \SI{91}{\percent} & $\checkmark$ \\
U05 & $\checkmark$ & 76 & 68 & $\checkmark$ & 358 & 299 & \SI{84}{\percent} & $\checkmark$ \\
U06 & $\times$ & 0 & 0 & $\times$ & 93 & 84 & \SI{90}{\percent} & $\checkmark$ \\
U07 & $\checkmark$ & 13 & 13 & $\checkmark$ & 132 & 124 & \SI{94}{\percent} & $\checkmark$ \\
U08 & $\times$ & 9 & 7 & $\times$ & 179 & 137 & \SI{76}{\percent} & $\checkmark$ \\
U09 & $\times$ & 0 & 0 & $\times$ & 301 & 291 & \SI{97}{\percent} & $\checkmark$ \\
U10 & $\times$ & 0 & 0 & $\times$ & 201 & 165 & \SI{82}{\percent} & $\checkmark$ \\
U11 & $\times$ & 0 & 0 & $\times$ & 216 & 155 & \SI{72}{\percent} & $\checkmark$ \\
U12 & $\checkmark$ & 37 & 30 & $\checkmark$ & 309 & 199 & \SI{64}{\percent} & $\checkmark$ \\
U13 & $\times$ & 4 & 4 & $\times$ & 202 & 115 & \SI{57}{\percent} & $\times$ \\
U14 & $\times$ & 5 & 3 & $\times$ & 420 & 398 & \SI{95}{\percent} & $\checkmark$ \\
U15 & $\checkmark$ & 30 & 23 & $\checkmark$ & 352 & 308 & \SI{88}{\percent} & $\checkmark$ \\
U16 & $\checkmark$ & 16 & 16 & $\checkmark$ & 346 & 222 & \SI{64}{\percent} & $\checkmark$ \\
U17 & $\checkmark$ & 21 & 21 & $\checkmark$ & 264 & 248 & \SI{94}{\percent} & $\checkmark$ \\
U18 & $\times$ & 4 & 4 & $\times$ & 210 & 188 & \SI{89}{\percent} & $\checkmark$ \\
U22 & $\times$ & 3 & 3 & $\times$ & 3 & 3 & \SI{89}{\percent} & $\checkmark$ \\
\hline
\end{tabular}

\end{table}
